\documentclass[11pt]{article}
\usepackage[margin=1.1in]{geometry}
\usepackage{amsmath,amssymb,amsthm}
\usepackage{listings}
\usepackage{xcolor}
\usepackage[hidelinks]{hyperref}
\usepackage{booktabs}

\newtheorem{theorem}{Theorem}
\newtheorem{lemma}{Lemma}
\newtheorem{proposition}{Proposition}
\newtheorem{corollary}{Corollary}
\theoremstyle{definition}
\newtheorem{definition}{Definition}
\newtheorem{remark}{Remark}
\newtheorem{example}{Example}

\lstset{
  basicstyle=\ttfamily\small,
  keywordstyle=\bfseries,
  columns=fullflexible,
  keepspaces=true,
  language=ML,
  morekeywords={val,type,of,let,match,with},
}

\newcommand{\bound}[1]{\mathrm{B}\,#1}
\newcommand{\ax}[2]{\mathrm{Axis}(#1,#2)}
\newcommand{\bc}[1]{\mathrm{Bcast}(#1)}
\newcommand{\grp}[1]{\mathrm{Group}[#1]}
\newcommand{\shape}{\mathrm{shape}}
\newcommand{\size}{\mathrm{size}}
\newcommand{\cosize}{\mathrm{cosize}}
\newcommand{\eval}{\mathrm{eval}}
\newcommand{\canon}[2]{\langle #1 \rangle_{#2}}
\newcommand{\logical}{\textrm{logical}}
\newcommand{\physical}{\textrm{physical}}
\newcommand{\tv}{\textrm{thread-value}}
\newcommand{\compose}{\mathrm{compose}}
\newcommand{\divideop}{\mathrm{divide}}
\newcommand{\repeatop}{\mathrm{repeat}}
\newcommand{\interleave}{\mathrm{interleave}}
\newcommand{\broadcast}{\mathrm{broadcast}}
\newcommand{\coalesce}{\mathrm{coalesce}}
\newcommand{\restrict}{\mathrm{restrict}}
\newcommand{\N}{\mathbb{N}}
\newcommand{\Z}{\mathbb{Z}}

\title{A Simplified Layout Algebra for GPU Kernels}
\author{Nan\thanks{\texttt{nan@silares.com}. \emph{TODO: full name and affiliation.}}}
\date{Draft of \today}

\begin{document}
\maketitle

\begin{abstract}
The layout algebra behind NVIDIA's CuTe library describes how tensor
tiles map onto threads and memory. A layout is a shape with a stride for each
axis, so the address of a coordinate is the sum of its indices times their
strides. As formalized by Shah and by Cecka, its
central operations, such as composition, complement, division, product, are
guarded by at least six interacting admissibility conditions. We show that this
complexity is an artifact of a representational choice, not of the domain.

We present a layout algebra in which every layout names the space its inputs
and outputs live in: logical indices, physical indices, or thread-value
indices. A layout into logical space must be a bijection from its $N$
coordinates onto $[0,N)$. A layout into physical space must be alias-free:
different coordinates land on different addresses. The exception is broadcast,
where many threads read one value on purpose, and the layout has to mark where
that happens. Nothing maps out of a physical index.

Composing $f$ with $g$ requires only that $f$'s range be exactly $g$'s
flattened set of coordinates. To emit the index arithmetic, naively, we can
always compute $f$, split that by division and remainder into $g$'s
coordinates, and feed it into $g$. Often the composite can be simplified to a
layout as well. We show that if it can be written this way the form is unique, and we emit that simplified
form.

A warp's or a thread's slice is a \emph{restriction}. A restricted layout is no
longer composable, which forces the user to settle the whole view of a tile
before taking slices.

We validate the algebra by showing that CuTe's operations, such as composition,
complement, and division, can be reproduced in it. Furthermore, we provide an
implementation in OCaml, and use it to rederive the atoms and usage patterns of
CuTe. 

The implementation is roughly 900 lines of code.
\end{abstract}

\section{Introduction}\label{sec:intro}

For GPU kernels, it is important to have a scheme for describing how to move
tiles of tensors between global memory, shared memory and registers. The de
facto vocabulary for such schemes is the CuTe layout algebra underlying
NVIDIA's CUTLASS library~\cite{cutedocs,cecka}. The most fundamental concept
is the layout, which we shall reuse heavily.

\begin{definition}[Layout]\label{def:layout}
A \emph{shape} is a positive integer $n$, or a tuple $(S_1,\dots,S_k)$ of
shapes; its \emph{size} is $\size n = n$ and
$\size (S_1,\dots,S_k) = \prod_i \size S_i$. A \emph{coordinate} of $S$ is
an integer $i \in [0,n)$, respectively a tuple $(c_1,\dots,c_k)$ of
coordinates of the $S_i$; write $\mathrm{Coord}(S)$ for the set of them. A
\emph{layout} $S{:}D$ is a shape $S$ together with a \emph{stride} $D$ of the
same structure: an integer where $S$ is a positive integer, and a tuple
$(D_1,\dots,D_k)$ where $S$ is $(S_1,\dots,S_k)$. Its index function
$\mathrm{Coord}(S) \to \N$ is defined by recursion on the shape,
\[
(n{:}d)(i) = i\,d,
\qquad
\big((S_1,\dots,S_k){:}(D_1,\dots,D_k)\big)(c_1,\dots,c_k)
= \sum_{j=1}^{k} (S_j{:}D_j)(c_j).
\] Flattening is a bijection $\canon{\cdot}{S} : \mathrm{Coord}(S) \to
[0,\size S)$, taken row-major, so that the last component varies fastest, with
inverse $\mathrm{unflatten}_S$. A coordinate may therefore be given as a single
integer, and we say that integer \emph{is} the coordinate. The \emph{cosize} of the layout is the least $N$ with its
image inside $[0,N)$, which is one plus the largest value it takes. The layout
is \emph{dense} when its image is an initial segment $[0,N)$ of the integers,
and a \emph{dense bijection} when it is moreover injective, in which case
$N = \size S$ and every value below $N$ is taken exactly once. A
\emph{mode} is one entry of the shape together with the matching entry of the
stride. Two layouts \emph{concatenate}: for $A = S{:}D$ and $B = S'{:}D'$, the
layout $(A,B)$ is $(S,S'){:}(D,D')$, and it sends $(c,c')$ to $A(c) + B(c')$.
\end{definition}

Layouts are fundamentally the correct concept here, because physical memory is
linearly addressed while the logical space is multidimensional and dense as we
imagine it. They are a natural way to abstract and describe those mappings.

The algebra has four operations. Three of them, composition, complement and
logical division, are formalized by Shah~\cite{shah}; the fourth, logical
product, we take from the CUTLASS implementation. The \emph{complement} of a
layout
$B = (N_0,\dots,N_\alpha){:}(d_0,\dots,d_\alpha)$, sorted by stride, with
respect to an integer $M$, is
\[
B^{*}_{M} \;=\;
\Big(d_0,\ \tfrac{d_1}{N_0 d_0},\ \dots,\ \tfrac{M}{N_\alpha d_\alpha}\Big)
\;{:}\;
\Big(1,\ N_0 d_0,\ \dots,\ N_\alpha d_\alpha\Big),
\]
chosen so that the concatenation $(B, B^{*}_{M})$ has size $M$ and its layout
function is a bijection of $[0,M)$ onto itself. \emph{Logical division} is
$A \mathbin{/} B = A \circ (B, B^{*}_{\size A})$, which reindexes $A$ so that
one mode runs inside the tile $B$ selects and the other runs across copies of
it. \emph{Logical product} is $A \otimes B = (A, A^{*}_{M} \circ B)$, where
$M = \size A \cdot \cosize B$, the smallest $M$ for which the composition is
defined. Its
first mode is one copy of $A$, and its second walks $B$ through the addresses
$A$ does not occupy, so the result is $\size B$ copies of $A$ placed in the
pattern $B$ names.

These operations are critical for constructing the layouts one usually wants.
There is, however, a lot of complexity underneath them. Layouts are in general
non-injective and non-surjective onto $[0,M)$, and CuTe always wants to write
the composite as a layout again. This causes at least six condition checks. Two are in the complement, whose
shape entries are quotients and so must divide: $N_{i-1}d_{i-1} \mid d_i$ for
every $i$, and $N_\alpha d_\alpha \mid M$. The others are in composition.
Write the flattened shape of $A$ as $M_0 \cdots M_\alpha$. It is \emph{left divisible} by $r$
when $r = M_0 \cdots M_{i-1} \cdot c$ for some $i$, where $c$ is smaller than
$M_i$ and divides it, and \emph{weakly left divisible} when the same holds
except that $c$ need not divide $M_i$. For each mode $(N){:}(r)$ of $B$ this is
three checks, since the prefix must divide $r$, the leftover $c$ must divide
$M_i$, and what remains after $r$ must be weakly left divisible by $N$. One
more is required across modes, that the intervals on which they act be pairwise
disjoint~\cite{shah}. None of these conditions is visible in the CuTe API,
yet all of them are visible failures.

We try to resolve those complexities by restricting what a layout can be.

\paragraph{Explicit index spaces.} The input and the output of a layout each
represent one of three things: a logical index into a tile, a thread-value
index, or a physical offset in memory. Therefore a layout should declare which of the three its
input and its output are, and we allow only the maps worth having between them.
Logical and thread-value indices map into one another in either direction.
Logical indices map into physical offsets. Nothing maps out of a physical
offset. By design, a thread reaches memory by composing its map into logical
space with the storage map.

\paragraph{Requirements that follow from the space.} We allow those maps and
not others because the three spaces do not carry the same obligations. Logical
and thread-value indices have to account for their contents losslessly, so a
map into either must be a dense bijection. Composition can only happen when one
layout's image and the other's domain agree completely, which is natural
because a physical offset is forbidden as a domain. A physical offset carries
far fewer requirements, since gaps in memory are common and sharing one address
between threads, to read the same element, is often what is wanted, so a map
into physical need only be alias-free, meaning that distinct coordinates reach
distinct addresses except where the layout marks a broadcast.

\paragraph{Composition and simplification.} The composite of two layouts is not
a layout in general. But we can decide if it can be written as a layout, and
when it is a layout we emit one weighted sum of the coordinates. When it is
not, we emit the composition literally, at the cost of a longer expression and
more work at run time.

\section{Our construction of layouts}\label{sec:layouts}

We build on top of Definition~\ref{def:layout}. First we define the spaces.

\begin{definition}[Spaces]\label{def:spaces}
Each side of a layout carries a space from $\{\logical, \physical, \tv\}$,
written $A \to B$. A layout may be annotated $A \to B$ for
$A, B \in \{\logical, \tv\}$, or $\logical \to \physical$, and no other pair is
constructible. Composition takes $A \to \logical$ and $\logical \to C$ to
$A \to C$, so $\tv \to \physical$ exists only as a composite.
\end{definition}

\begin{definition}[Replication]\label{def:linear}
Extend Definition~\ref{def:layout} by allowing the stride against a positive
integer $n$ to be the mark $\bot$ in place of an integer, with
\[
(n{:}\bot)(i) = 0 \quad\text{for every } i \in [0,n).
\]
We denote $\ax{n}{s}$ for $n{:}s$ and $\bc{n}$ for $n{:}\bot$.
\end{definition}

A zero stride is rejected rather than read as replication. This is because
replication is legal for reads and not for writes, since simultaneous loads of
one address are a broadcast while simultaneous stores are a race. It is better
not to use $0$, so that the two can never be mixed.

\begin{definition}[Swizzled layout]\label{def:swizzle}
A \emph{swizzle} is a bijection $\physical \to \physical$ of the form
\[
\sigma(x) \;=\; x \oplus \big(((x \gg \mathit{src}) \wedge (2^{b}-1)) \ll
\mathit{dst}\big),
\]
its two fields $[\mathit{src},\mathit{src}{+}b)$ and
$[\mathit{dst},\mathit{dst}{+}b)$ disjoint. A layout $A \to \physical$
followed by a swizzle is a \emph{swizzled layout}, again $A \to \physical$.
\end{definition}

A swizzle runs from $\physical$ to $\physical$, so it can only ever be the last
step, and being a bijection it sends distinct addresses to distinct addresses.
It therefore leaves the conditions below untouched, and matters only in the
arithmetic that is finally emitted.

\begin{definition}[Composition]\label{def:compose}
Let $f : A \to \logical$ and let $g : \logical \to C$ have domain shape $S_g$.
The \emph{composite} $g \circ f : A \to C$ is defined when $f$ is a dense
bijection onto $[0, \size S_g)$, and is the map
\[
g \circ f \;=\; c \mapsto g\big(\mathrm{unflatten}_{S_g}(f(c))\big),
\]
whose domain shape is that of $f$.
\end{definition}

\begin{definition}[Declared replicas]\label{def:replica}
For a layout $L$ with domain shape $S$, define a relation $\sim_L$ on
$\mathrm{Coord}(S)$:
\begin{itemize}
\item if $L$ is a layout, $c \sim_L c'$ when $c$ and $c'$ agree at every
position whose stride is an integer, that is, they differ only at positions
marked $\bot$;
\item if $L = g \circ f$ with $g$ of domain shape $S_g$, then $c \sim_L c'$ when
$\mathrm{unflatten}_{S_g}(f(c)) \sim_g \mathrm{unflatten}_{S_g}(f(c'))$.
\end{itemize}
\end{definition}

\begin{definition}[Validity]\label{def:validity}\label{sec:validity}
Let $L : A \to B$ be a layout with domain shape $S$. Then $L$ is \emph{valid}
when
\[
\begin{cases}
L \text{ is a dense bijection}
  & \text{if } B \in \{\logical, \tv\},\\[2pt]
L(c) = L(c') \implies c \sim_L c'
  & \text{if } B = \physical.
\end{cases}
\]
It is \emph{write-valid} when $B = \physical$ and in addition
\[
L(c) = L(c') \implies c = c'.
\]
Only a write-valid layout may be the target of a store.
\end{definition}

Whether a layout is a dense bijection can be decided without enumeration.

\begin{lemma}[Dense bijections]\label{lem:dense}
Let $L$ be a layout with no marked positions, its positions carrying sizes
$n_1,\dots,n_k$ and strides $s_1,\dots,s_k$ with every $n_i \ge 2$, and put
$N = \prod_i n_i$. Then $L$ is a dense bijection onto $[0,N)$ if and only if,
after sorting the positions by stride, $s_{(1)} = 1$ and
$s_{(j+1)} = s_{(j)}\,n_{(j)}$ for every $j$: the sorted strides are the
running products of the sorted sizes.
\end{lemma}

\begin{proof}
($\Leftarrow$) With running-product strides the map is exactly mixed-radix
representation with radices $n_{(1)},\dots,n_{(k)}$, a bijection onto $[0,N)$.
Positions of size $1$ contribute nothing and are dropped beforehand.

($\Rightarrow$) Sort the positions by stride and write $N_j = \prod_{i \le j} n_{(i)}$
(so $N_0 = 1$). \emph{Claim:} $s_{(j)} \le N_{j-1}$ for every $j$. Every image
value below $s_{(j)}$ comes from a coordinate supported on positions
$1,\dots,j{-}1$, since a nonzero index at any later position contributes at
least $s_{(j)}$ and no contribution is negative; those positions produce at
most $N_{j-1}$ distinct values, so surjectivity onto $[0,s_{(j)})$ forces
$s_{(j)} \le N_{j-1}$. Consequently
\[
\textstyle\sum_j (n_{(j)}-1)\,s_{(j)} \;\le\; \sum_j (n_{(j)}-1)\,N_{j-1}
\;=\; N - 1,
\]
the last equality by telescoping. But the left side is the maximum of the
image, and surjectivity onto $[0,N)$ requires the maximum to be exactly
$N-1$. The two sums therefore agree, and since each term on the left is at most
the corresponding term on the right, they agree termwise; as
$n_{(j)} - 1 \ge 1$ this gives $s_{(j)} = N_{j-1}$ for every $j$, which is the
running-product condition.
\end{proof}

\section{Operations}\label{sec:ops}

The operations below transform layouts, and their behaviour on a swizzle is
stated where it matters. An \emph{arrangement} is a layout required to be a
dense bijection, checked by Lemma~\ref{lem:dense} at the call; its strides
carry order and nothing else. An arrangement cannot express padding, because a
gap in it would be a whole footprint wide. Padding lives in the strides of the
tile being copied, and the tile's cosize then spaces the copies.

\begin{description}
\item[$\mathrm{canonical}(S)$] the row-major layout of $S$: last axis stride
1, running products leftward. Dense by Lemma~\ref{lem:dense}.

\item[$\mathrm{split}(\mathit{by}{:}S)$] regroups one position
$\ax{n}{s}$ with $\size S = n$ into the structure of $S$, giving each new
position the stride $s$ times its running product. The index $i$ is read as the
$S$-coordinate of $i$, so the map is unchanged. A marked position splits into
marked positions.

\item[$\divideop(\mathit{by}{:}S)$] \emph{the partition} (CuTe's
zipped\_divide with a shape tiler). Per axis, size $n_i$ splits into
$(n_i/k_i, k_i)$ (divisibility required); the pieces regroup as
$\grp{\text{tile axes}; \text{rest axes}}$ with coordinate
$((\text{tile}),(\text{which tile}))$. It is the relabeling
$\rho((t_i),(b_i)) = (b_i k_i + t_i)_i$, a bijection, with
$\divideop\, t = t \circ \rho$. Every tile is kept, so the rest needs no
complement operation. Selecting one tile is restriction
(Section~\ref{sec:restrict}), afterwards.

\item[$\repeatop(\mathit{by})$] copies side by side:
$\grp{\mathit{by} \times \cosize t;\; t}$, the arrangement's strides scaled to
footprint units. A dense arrangement over a dense tile is dense, since
$\cosize t = \size t$ there and the two digit systems concatenate.
On a swizzled layout, $\repeatop$ is accepted iff both swizzle fields lie
below the (power-of-two) footprint: multiples of the footprint then commute
with the swizzle, so each copy is independently swizzled (the multi-stage
shared-memory buffer); otherwise it rejects---the copies would contaminate one
another's swizzle.

\item[$\interleave(\mathit{by})$] copies dealt element-wise (CuTe's raked
product), the mirror image: $\grp{\mathit{by};\; t \times \size(\mathit{by})}$.
Blocked versus raked is the coalescing distinction: with threads as the
arrangement, $\interleave$ gives consecutive lanes consecutive addresses.
No meaning is defined for interleaving a swizzled layout (no kernel pattern
rakes swizzled tiles); it rejects.

\item[$\broadcast(\mathit{by}{:}S)$] copies at the \emph{same} addresses:
$\grp{\mathrm{Bcast}\text{-tree of } S;\; t}$. The dual of $\repeatop$; the
arrangement is a shape because aliased copies have no positions to give.
Read-legal, write-illegal by the $\mathrm{Bcast}$ rule.

\item[$\coalesce$] flattens nesting, drops positions of size $1$, merges
$\ax{a}{s'}$ into an inner neighbour $\ax{b}{s}$ when $s' = s\,b$, and merges
adjacent marked positions. The map is unchanged. It is an optional normal form
here, never a precondition of anything.

\item[$\mathrm{inverse}$] of a dense bijection, defined on layouts: the index
$j$ maps to the coordinate sent to $j$, whose position $i$ carries index
$(j / s_i) \bmod n_i$. The spaces swap, so the inverse of a thread--value map
$\tv \to \logical$ is the \emph{ownership map} $\logical \to \tv$, which
element is held by which thread and value. Its use is \emph{retiling}: two
thread-partitionings of one tile, a copy's fragment and an MMA's operand, are
related by $\mathrm{inverse}(\mathit{mma}_{TV}) \circ \mathit{copy}_{TV}$,
carrying a copy coordinate to an element to an MMA coordinate. That
correspondence is a bijection and need not fix the lane: with 32 lanes holding
two values each, $\mathit{copy} = 2\,\mathit{lane} + v$ against
$\mathit{target} = \mathit{lane} + 32\,v$ sends 62 of the 64 coordinates to a
different lane. A composite is a bijection too, so its inverse function exists, but that
function need not be a layout. The layout $(3,2){:}(1,3)$ sends the six indices
to $0,3,1,4,2,5$, and composed with itself it sends them to $0,4,3,2,1,5$: that
map is its own inverse and is a layout over no refinement of six. An inverse
may instead be kept as a composite, with the factors inverted in reverse order.
That needs every factor invertible, so it applies only to chains that stay left
of $\physical$; storage is alias-free rather than injective and may leave gaps,
so a composite ending in memory has no inverse at all.
\end{description}

\begin{proposition}[Closure]\label{prop:closure}
$\mathrm{split}$, $\divideop$, and $\coalesce$ preserve the function up to the
stated relabeling, hence preserve density and injectivity. $\repeatop$ and
$\interleave$ with dense arguments preserve density; with an injective tile
they preserve injectivity. $\broadcast$ preserves the image and multiplies
multiplicity uniformly.
\end{proposition}
\begin{proof}
Split and coalesce regroup digits in place, leaving the index unchanged, and
divide composes per-axis splits with the permutation $\rho$. For $\repeatop$
the images are
$\{a \cdot \cosize t + v : a \in \mathrm{Im}(\mathit{by}), v \in
\mathrm{Im}(t)\}$; density of $\mathit{by}$ makes $a$ range over
$[0,\size \mathit{by})$ and, for dense $t$, $\cosize t = \size t$, so the two
digit systems concatenate. Injectivity survives a tile with gaps because
distinct copies are separated by at least the footprint. $\interleave$ is
symmetric.
\end{proof}

\section{Composition}\label{sec:compose}

Composition is the only place two layouts meet, and the type does most of the
work:
\[
\compose : (\alpha, \logical)\ \mathrm{Layout} \to (\logical, \gamma)\
\mathrm{Layout} \to (\alpha, \gamma)\ \mathrm{Layout}.
\]
Physical cannot be an intermediate because no signature accepts it on the
left, and two thread--value maps cannot be composed with each other because
they share a codomain, not an interface. The composition chain is therefore
short and terminal: $(\tv \to \logical)$ maps compose through zero or more
$(\logical \to \logical)$ relabelings into one $(\logical \to \physical)$
storage map, after which nothing composes.

The condition on $f$ is decided by Lemma~\ref{lem:dense} when $f$ is a layout
and by enumeration when it is a composite, which happens in chains. It rejects
a marked position on the left, since replication is not injective. Nothing else
is checked.

CuTe's \texttt{composition} of two layouts must return a layout, so it has to
align the operands' digit structures and succeeds only when their shapes and
strides divide one another suitably. It carries a second representation,
\texttt{ComposedLayout}, for the cases where that fails, and that
representation has a reduced contract: no \texttt{stride()}, alignment queries
pinned to $1$, a zero-offset requirement when it is the right operand of a
composition with an ordinary layout, and structural operations delegated to one
factor. Two separable things are at work here. The spaces and their requirements decide
which maps may be built at all, which is what CuTe leaves to convention.
Keeping the operand pair is what removes the divisibility conditions, and
\texttt{ComposedLayout} shows that CuTe can keep a pair too. The combination is
what differs: one representation with one contract, whether or not the operands
would have fused, and the arithmetic recovered afterwards by
Theorem~\ref{thm:decide}, exactly when it exists.

\begin{example}[A composite that is not a layout]\label{ex:nonstrided}
Let $f$ be the relabel $(a,b) \mapsto 3a{+}b$ of shape $(2,3)$ and let $g$ be
the storage $(3,2){:}(1,3)$. Both are valid, and the composite sends the
size-3 axis at $a = 0$ to the addresses $0,3,1$. No layout has this map, since
along one axis a layout runs in an arithmetic progression and a prime-sized
axis cannot be split, so a composition that must return a layout has to reject
this pair. Here the composite exists and is a bijection onto $[0,6)$, emitted
as $\lfloor x/2 \rfloor + 3(x \bmod 2)$ for $x = 3a{+}b$, and
Theorem~\ref{thm:decide} proves that no layout describes it.
\end{example}

\subsection{Emitting a composite}\label{sec:emit}

The emitter writes Definition~\ref{def:compose} out as an integer expression
over the composite's coordinate variables:
\[
E \;::=\; k \;\mid\; c \;\mid\; \textstyle\sum_i a_i E_i + k \;\mid\;
\lfloor E / w \rfloor \;\mid\; E \bmod r \;\mid\; E \oplus E \;\mid\;
\mathrm{let}\ x = E\ \mathrm{in}\ E,
\]
with $k, a_i, w, r$ compile-time constants and each coordinate variable $c$
ranging over $[0, n_c)$, the size of its position. A layout with strides $s_i$
writes $\sum_i s_i c_i$, and a marked position contributes nothing. A composite
binds the index $f$ computes and reads $g$'s digits off it: digit $j$,
of radix $n_j$ and weight $w_j$ (the product of the radices to its right), is
$\lfloor x / w_j \rfloor \bmod n_j$, without the modulus on the leading digit
since composability gives $x < \size(g)$. A swizzle likewise binds
the address it is given and exclusive-ors one shifted field of it. Binding is
what keeps this small: a stage index is computed once however many digits are
taken from it. This step cannot fail; it is the pair, written out. For the
running example---the thread--value map $(cg,(r,v)) \mapsto 4cg + 16r + v$
onto column-major $(8,16)$ storage---it writes
\begin{center}\small
\texttt{let x1 = 4*c0 + 16*c1\_0 + c1\_1 in (x1 / 16) + 8*(x1 \% 16)}.
\end{center}

When the composite is a layout, its address is a single weighted sum, and this
expression hides that. Recovering it looks like a simplification problem, and
we first treated it as one, with a bounds-aware rewrite system of the kind
Halide~\cite{halide}, TVM~\cite{tvm} and MLIR's affine dialect~\cite{mlir}
carry. That was the wrong tool. It missed layouts that exist, and where none
exists it made the code worse, since exposing cancellations means inlining the
bound index into every digit and splitting the digits apart: on random operands
that inflated the arithmetic by 45\% over the expression above, in 400 of 628
cases. The next subsection decides the question instead.


\subsection{Deciding whether a composite is a layout}\label{sec:decide}

A composite is a function on a finite box, and at compile time every one of its
values is available. Whether it is a layout, over a possibly finer shape, is
therefore decidable by evaluation rather than by searching for a
representation. The procedure below costs $O(\size)$ evaluations, which is what
the checks of Section~\ref{sec:validity} already spend, plus work linear in
each axis.

Fix a map $f : B \to \N$ on the box $B = \prod_{i=1}^{m} [0, n_i)$, and write
$e_i$ for the unit coordinate on axis $i$ and $\mathbf{0}$ for the origin.

\begin{definition}[Refinement; layout over a refinement]\label{def:refine}
A \emph{refinement} of axis $i$ is a strictly increasing chain of divisors
$1 = w_1 \mid w_2 \mid \cdots \mid w_k \mid n_i$ with $w_k < n_i$, presenting
$c_i$ by its digits $d_j(c_i) = \lfloor c_i / w_j \rfloor \bmod r_j$ with
$r_j = w_{j+1}/w_j$ and $w_{k+1} = n_i$. Starting at $1$ and ending below $n_i$
is what makes the digits \emph{cover} the axis: every radix is at least $2$ and
$\prod_j r_j = n_i$, so $(d_1,\ldots,d_k)$ is a mixed-radix presentation of
$c_i$. For $n_i = 1$ the chain is empty and there are no digits. A
\emph{layout over the refinement} is
$f(c) = k + \sum_{i,j} s_{ij}\, d_j(c_i)$ with integer $k$ and $s_{ij}$,
where a stride may be zero. The trivial refinement, one digit of radix $n_i$
per axis, gives the forms that are affine in the coordinates themselves.
\end{definition}

Being a layout over a refinement of the domain shape is what CuTe's
composition must produce, and what its admissibility
conditions are conditions for.

\begin{lemma}[Separability]\label{lem:separable}
$f$ is a layout over some refinement if and only if
\begin{equation}\label{eq:sep}
f(c) = f(\mathbf{0}) + \sum_{i=1}^{m} g_i(c_i), \qquad
g_i(v) = f(v\,e_i) - f(\mathbf{0}),
\end{equation}
and each $g_i$ is a layout over some refinement of $[0, n_i)$.
\end{lemma}
\begin{proof}
($\Rightarrow$) Every digit is a function of one coordinate, so such a layout
is $f(c) = k + \sum_i h_i(c_i)$ with $h_i = \sum_j s_{ij} d_j$. Evaluating at
$\mathbf{0}$ and at $v\,e_i$ gives $h_i(v) - h_i(0) = g_i(v)$, so $g_i$ is
$h_i$ shifted by a constant---a layout over the same refinement---and
\eqref{eq:sep} holds. ($\Leftarrow$) Sum the per-axis forms.
\end{proof}

Separability is one pass over the box. It leaves $m$ independent
one-dimensional problems, and each is settled in time linear in its axis.

\begin{lemma}[Floor form]\label{lem:floor}
For $w \mid w'$, $\lfloor v/w \rfloor \bmod (w'/w) =
\lfloor v/w \rfloor - (w'/w)\lfloor v/w' \rfloor$. Hence $g$ is a layout
over a chain $w_1 \mid \cdots \mid w_k \mid n$ if and only if
$g(v) = \sum_j a_j \lfloor v/w_j \rfloor$ for integers $a_j$ on the same chain,
and the two coefficient families determine each other by
$a_j = s_j - r_{j-1}s_{j-1}$ and $s_j = g(w_j)$.
\end{lemma}
\begin{proof}
The identity is the definition of the remainder, valid because $w \mid w'$.
Substituting it into $\sum_j s_j d_j$ and collecting the terms in
$\lfloor v/w_j \rfloor$ gives $a_j = s_j - r_{j-1}s_{j-1}$ with $s_0 = 0$;
the relation is triangular with unit diagonal, hence invertible over $\Z$.
Evaluating $\sum_i a_i \lfloor w_j/w_i \rfloor$ at $v = w_j$ kills every term
with $w_i > w_j$ and leaves $\sum_{i \le j} a_i (w_j/w_i)$, which the same
recurrence identifies with $s_j$.
\end{proof}

\begin{theorem}[Linear recognition]\label{thm:decide}
Let $g : [0,n) \to \Z$ with $g(0) = 0$ and let $h(t) = g(t) - g(t-1)$ for
$1 \le t < n$. If $n = 1$, return the empty refinement. Otherwise begin
the refinement with $w_1 = 1$, set $w_{\mathrm{last}} = 1$, and process
$t = 1, 2, \ldots, n-1$ in order: whenever $h(t) \ne 0$, reject unless
$t \mid n$ and $w_{\mathrm{last}} \mid t$; otherwise record the coefficient
$a_t = h(t)$, append $t$ to the refinement if $t > 1$, set
$w_{\mathrm{last}} = t$, and subtract $a_t$ from $h$ at every multiple of
$t$. Then $g$ is a layout over some refinement if and
only if the scan completes without rejecting; the reported weights are
\emph{necessary} boundaries of any such refinement, so together with $w_1 = 1$
they are the weights of the coarsest one; the strides are $s_j = g(w_j)$; and
the whole scan runs in $O(n)$ time.
\end{theorem}
\begin{proof}
By Lemma~\ref{lem:floor} such a layout is equivalent to
$g = \sum_j a_j \lfloor \cdot/w_j \rfloor$ over a divisor chain ending at $n$.
First differences of $\lfloor v/w \rfloor$ are the indicator of the multiples
of $w$, so $h(t) = \sum_j a_j [\,w_j \mid t\,]$. Induct on $t$: when the scan
reaches $t$, the contributions of all weights below $t$ have been subtracted,
so the residual $h(t)$ equals $a_t$, the coefficient at weight $t$, which is
therefore forced---and nonzero exactly when $t$ is a boundary. Necessity of
the divisibility tests is immediate: the weights of a refinement divide $n$
and form a chain. Conversely, if the scan completes, the reported weights are
a chain of divisors of $n$ and the recovered $a_j$ reproduce $g$ by
construction, so Lemma~\ref{lem:floor} converts them to strides $s_j = g(w_j)$.
Since each $a_t$ is forced, a boundary with $a_t \ne 0$ appears in every
refinement admitting a form, which is the claimed necessity. For the running
time, each accepted weight is a strictly larger multiple of its predecessor,
so accepted weights at least double; the subtraction loops therefore perform
fewer than $n + n/2 + n/4 + \cdots < 2n$ updates, alongside the single scan.
Finally, $w_1 = 1$ belongs to the reported refinement whatever $a_1$ is,
because every refinement carries it (Definition~\ref{def:refine}); a zero
coefficient there means a zero stride, not an absent digit. For
$g(v) = \lfloor v/2 \rfloor$ on $n = 4$ the differences are $0,1,0$ and only
weight $2$ is reported, but the refinement is $(2,2)$ with strides $(0,1)$, and
omitting its first digit would leave radices whose product is not $n$. A
constant $g$---necessarily $0$, since $g(0)=0$---gives the trivial refinement
at stride $0$.
\end{proof}

The characterization these two lemmas state---digits as differences of floor
terms, first differences as indicators of divisibility, and the coefficients
they force---is also the machinery of Appendix~A.2.1 of Axe~\cite{axe}, whose
Lemmas~A.3--A.6 and Theorem~A.8 use it to prove a normalized sharded
representation unique by reconstructing it from the function it induces. We
claim no novelty for it. What Theorem~\ref{thm:decide} adds is the question
asked of it: the input is an arbitrary map on a finite box rather than one
assumed representable, so existence is \emph{decided} and non-representable
maps are rejected; the extraction is an incremental sieve rather than M\"obius
inversion over the divisor lattice, which is what gives the linear bound; and a
rejection is not a dead end, because the emitter has the pipeline of
Section~\ref{sec:emit} to fall back on.

No factorization and no search over orderings appears anywhere, since the
boundaries announce themselves. Example~\ref{ex:nonstrided} is settled by the
scan rather than argued: its size-3 axis has differences $3, -2$, forcing a
weight at $1$ and again at $2$, and $2 \nmid 3$. Because the procedure reads
values, a swizzle is no obstacle to it; the layout $2c$ under $\{1;1;0\}$ emits
as $(2c) \oplus (\lfloor 2c/2 \rfloor \bmod 2)$ and is decided to be $3c$. A
restriction is decided the same way over its free coordinates, with the pinned
contribution folded into the constant.

The emitter therefore has two outputs. If the map is a layout over some
refinement it emits that layout's arithmetic, over the coarsest refinement and
with the coefficients the map forces, so no form with fewer digits exists.
Otherwise it emits the expression of Section~\ref{sec:emit}.


\section{Restriction}\label{sec:restrict}

A kernel does not use a whole address map. Warp $w$ and thread $t$ each use a
slice of one. Fixing a coordinate cannot produce a layout, since the fixed
contribution is a position rather than a pattern, and it cannot be handed back
as a separate offset either, because $\sigma(64 + x) \ne 64 + \sigma(x)$
whenever the swizzle reads the fixed bits. So the fixed coordinates are kept
as data beside the layout.

\begin{definition}[Partial coordinate; restricted layout]
A \emph{partial coordinate} of shape $S$ is
$\pi ::= \mathrm{Free} \mid \mathrm{At}\,i \mid
\mathrm{Parts}[\pi_1,\dots,\pi_k]$, structurally matching $S$ with in-bounds
indices. A \emph{restricted layout} is a pair $(L, \pi)$ of an
\emph{unmodified} layout and a partial coordinate of its domain shape. Its
shape replaces fixed slots by $\bound 1$; its offset function
\emph{completes} the coordinate (fixed slots take index $0$ and contribute
their $\mathrm{At}$ index) and evaluates $L$.
\end{definition}

Nothing is folded or split, so restriction applies to every map, layout or
composite, swizzled or not, with no case analysis. Only a physical codomain may
be restricted: choosing which addresses this executor touches is an execution
decision, whereas density in logical space forbids discarding elements. Slicing
therefore happens after composition, where the warp and thread coordinates are
slots of the composite's domain. Restricting again takes the union of the two
partial coordinates, so slicing by warp and then by thread is associative.

Because layouts and restrictions are closed data, a slice needs no separate
treatment: the restricted map is decided in its own right over the free
coordinates, with the pinned contribution folded into the constant, and falls
back to the pipeline with those variables replaced by their constants. Warp~2's
slice of the composite above emits $r + 8v + 64$. The per-lane math is
generated from the layout, not tabulated.

\section{CuTe's operations in this algebra}\label{sec:cute}

We now express CuTe's central operations in this algebra and execute the
correspondence on NVIDIA's own worked examples~\cite{cutedocs}. (Convention
translation is mechanical: a CuTe layout's mode list reverses, recursively,
into our row-major tree; CuTe's $\mathrm{composition}(A,B)$ applies $B$ first
and is our $\compose$ with $f = B$, $g = A$.)

\paragraph{Composition.} All three documented examples are reproduced as
functions---our composite agrees pointwise with the strided layout CuTe
prints, including the mode-splitting example
$\mathrm{composition}\big((6,2){:}(8,2),\,(4,3){:}(3,1)\big) =
\big((2,2),3\big){:}\big((24,2),8\big)$. Every published example has a dense
right operand, so all of them compose here.

\paragraph{Complement.} There is no complement operator here. $\divideop$
carries the rest by construction, and the density check decides complement's
defining property. Both directions of
$\mathrm{complement}(4{:}1, 24) = 6{:}4$ are witnessed by the single
$\divideop$ of $\mathrm{canonical}(24)$ by $\bound 4$, since tile and rest are
each other's complements. The strided case $\mathrm{complement}(4{:}2,24) =
(2,3){:}(1,8)$ comes from dividing the $(12,2)$ reshape. All six documented
complements are checked to make $(B, B^{*})$ dense onto $[0,24)$.

\paragraph{Division.} CuTe's strided-tiler example
$\mathrm{logical\_divide}\big((4,2,3){:}(2,1,8),\, 4{:}2\big)$ equals our
composition with the completed tiler $(B, \mathrm{complement}(B,24))$ on the
left, which is their division carried out as a composition with a dense
bijection, on their numbers.

\paragraph{Product, and one divergence.} Their $\mathrm{logical\_product}$
example is verified against their formula. On a tile with gaps the two differ:
CuTe's product places copies into the tile's gaps, while $\repeatop$ keeps
them, since padding is a property of the storage and no operation here
repurposes it.

\paragraph{Which conditions go, and which stay.} Two kinds of condition have to
be told apart. The divisibility conditions exist so that composition and
complement can return a layout: left divisibility, weak left divisibility,
disjoint intervals of definition, complement's two divisibility requirements,
and division's lemma that the pair it assembles is admissible in turn. Nothing
here fuses, so all of those go, and complement goes with them because division
carries its rest.

What stays is not empty, and none of it is about representation. Composition
requires the left operand to cover the right operand's domain exactly, which is
stronger than the containment ordinary function composition needs, and is
decided for layouts by Lemma~\ref{lem:dense}. A map into $\logical$ or $\tv$
must be a dense bijection, a map into $\physical$ must be alias-free, and a
store target must be injective. Three operations carry requirements of their
own: $\divideop$ needs each axis size divisible by its tile size, $\repeatop$
on a swizzled layout needs both swizzle fields below the footprint, and
$\interleave$ rejects a swizzled layout. These are obligations about what a map
means, not about whether two maps can be written as one.

\section{Coverage of CuTe usage}\label{sec:idioms}

The previous section covered CuTe's operators. This one covers the patterns
CUTLASS kernels are written with, and how each is expressed here. The last
column names the suite that executes the expression (Section~\ref{sec:eval});
a dash means it is expressible but no test exercises it.

\begin{center}\small
\begin{tabular}{@{}p{0.27\linewidth}p{0.36\linewidth}p{0.29\linewidth}@{}}
\toprule
\textbf{CuTe idiom} & \textbf{Here} & \textbf{Verified} \\
\midrule
\texttt{make\_layout(shape, stride)} & a layout & throughout \\
\texttt{composition(A, B)} & $\compose(B, A)$ (kept as the pair) & Suite 1, all three documented examples \\
\texttt{complement(B, M)} & the rest component of $\divideop$; its defining property decided by the density check & Suite 1, all six documented examples \\
\texttt{logical\_divide}, \texttt{zipped\_divide} & $\divideop$: $((\text{tile}),(\text{rest}))$, every tile kept & Suite 1 (strided-tiler example via the completed tiler); Suites 2--3 \\
\texttt{local\_tile} & $\divideop$, then $\restrict$ on the composite pinning the rest coordinate to this block's tile & Suite 2 (warp-tile derivation), Suite 3 \\
\texttt{local\_partition}, \texttt{partition\_A/B/C}, \texttt{tensor(\_, k)} & $\restrict$ on the composite: pin the executor's coordinates & Suite 3 (stage slice), warp-tile derivation \\
\texttt{logical\_product}, \texttt{blocked\_product}, \texttt{tile\_to\_shape} & $\repeatop$ with the order as a dense arrangement & Suite 2 (staged buffer; SM90 GMMA accumulator) \\
\texttt{raked\_product} & $\interleave$ & Suite 2 (blocked vs raked), Suite 3 (copy partition) \\
stride-0 modes (replication) & $\broadcast$ / marked positions, read-only by rule & Suite 2 (GMMA operand replication) \\
\texttt{Swizzle<B,M,S>} $\circ$ layout & $\mathrm{with\_swizzle}\ \{B; \mathit{src}; \mathit{dst}\}$ on a storage layout & Suite 2 (SW32/64/128 atoms, bank property) \\
multi-stage pipeline buffer & $\repeatop$ of the swizzled tile (legality: fields below the footprint) & Suite 2 (3-stage SW128), Suite 3 \\
\texttt{TiledCopy}: \texttt{partition\_S/D} with vector width & $\divideop$ by the vector, $\interleave$ over lanes (lanes fast $\Rightarrow$ coalesced) & Suite 3 (coalescing, vectorization asserted) \\
\texttt{TiledMMA}: fragment over a CTA tile & fragment atom $\interleave$d by the warp grid, $\compose$d onto the tile & Suite 2 (SM90 GMMA accumulator from the SM80 $C$ fragment), warp-tile derivation \\
\texttt{right\_inverse}, \texttt{retile\_S/D} & $\mathrm{inverse}$: the ownership map; retiling is the correspondence $\compose(\mathit{copy}_{TV}, \mathrm{inverse}(\mathit{mma}_{TV}))$ & \texttt{test\_layouts} (inverse law, retile correspondences) \\
\texttt{coalesce} & $\coalesce$, optional & \texttt{test\_layouts} \\
\texttt{group\_modes}, \texttt{select}, \texttt{flatten} & tree regrouping of $\mathrm{Group}$; $\coalesce$ & --- \\
\texttt{recast}, \texttt{upcast} & $\divideop$ by the vector width; the unit change is outside the algebra & --- \\
TMA box, multicast & $\divideop$ by the box shape; $\broadcast$ across CTAs & --- (box), Suite 2 (replication) \\
identity tensor + predication & masking, outside the layout algebra (Section~\ref{sec:scope}) & --- \\
\bottomrule
\end{tabular}
\end{center}

The middle column uses eight operations, namely $\compose$, $\divideop$,
$\restrict$, $\repeatop$, $\interleave$, $\broadcast$,
$\mathrm{with\_swizzle}$ and $\mathrm{inverse}$, together with $\coalesce$. No
row needs an operation absent from Section~\ref{sec:ops}.

Four rows are unexercised. Mode surgery and unit changes regroup a shape
without touching composition or tiling. Predication is masking, which
Section~\ref{sec:scope} places outside the algebra. The TMA box is $\divideop$
by the box shape, so it uses nothing the verified tiling rows do not; what is
untested there is the descriptor it feeds, which is not a layout question.
Composition, tiling and partitioning appear only in verified rows.

\section{Evaluation}\label{sec:eval}

The evaluation is four executable suites over an implementation of roughly
630 lines of OCaml for the algebra, 130 for the decision procedure and 150 for
the expression type it emits (1{,}250 with interfaces and their
documentation), and 1{,}700 lines of tests; every law stated in this paper is
asserted in them.

\paragraph{Transcription and derivation.} A layout enters a test suite one of
two ways. A \emph{transcription} copies a published layout and checks
properties of the data, which tests the representation. A \emph{derivation}
builds the layout from the operations and compares it against the
transcription, which tests the algebra. A missing operation is invisible to
transcription, because the transcriber performs it mentally, and $\divideop$
went unnoticed here for exactly that reason until a derivation needed it.
Hardware facts, the PTX fragment tables and the swizzle parameters, are the
only layouts we transcribe; everything built from them is derived.

\paragraph{Emission oracle.} Every map any suite builds, 178 of them, is
emitted both ways, as the expression of Section~\ref{sec:emit} and as the
decided form, and both are evaluated at every coordinate against the map's own
evaluation. A form affine in the coordinates is fitted to each by enumeration,
from the origin and the unit coordinates and then checked at every point, and
the harness asserts that whenever the fit succeeds the decision procedure
returns exactly that sum, never a different form and never ``none''. Of the
178, 147 emit as one weighted sum, 9 as digits of a single coordinate, and 22
as the expression.

\paragraph{Suite 4: random layouts.} Random operands of the two shapes
composition sees, a dense relabel from a random factorization and a storage map
of random axis order with random padding. 600 pairs, 300 three-stage chains and
300 swizzled composites, all checked at every coordinate under the same oracle.
Of the pairs, 327 are affine in their coordinates and every one is decided to
that sum, 524 are layouts over some refinement, and the decided forms use 686
divisions and moduli against the 1{,}026 of the emitted expression. Of the
chains, 119 are affine and all 119 are decided, which is the case the rewrite
system missed.

\paragraph{Suite 1: CuTe's central algebra} (Section~\ref{sec:cute}): three
composition examples, six complements, the strided-tiler division, the product
example, and the documented divergence.

\paragraph{Suite 2: hardware atoms and layouts derived from them.} The
transcribed atoms are checked against PTX: SM70 quadpair fragments, the SM80
\texttt{mma.m16n8k16} $A$, $B$ and $C$ fragments, the \texttt{ldmatrix}
$m8n8$ fragment, and the GMMA shared-memory swizzle atoms. The last are checked
against the CUTLASS source to be the family
$\{B; \mathit{src}{=}6; \mathit{dst}{=}3\}$ over $(8, 8\cdot 2^B)$
half-elements for $B \in \{0,1,2,3\}$, and the SW128 bank property, that the
unit index becomes $u \oplus r$, is asserted. Two composites are then derived
from those atoms rather than transcribed: a three-stage swizzled pipeline
buffer, and the SM90 GMMA $64{\times}128$ accumulator layout, built from the
SM80 \texttt{mma.m16n8k16} $C$ fragment by $\divideop(16,8)$, then
$\interleave$ by the canonical $(4,16)$ rest grid, then composition. It equals
NVIDIA's printed layout at all 8192 coordinates.

\paragraph{Suite 3: a GEMM mainloop from the operations alone.} The tile is
$(8,32)$ in halves, global leading dimension 40, 16-byte vectors, one copying
warp, and two shared-memory stages of the SW64 atom. The writer is
$\divideop(1,8)$ then $\interleave$ over 32 lanes. The reader is
$\divideop(8,8)$ then $\interleave$ of the \texttt{ldmatrix} atom over the four
matrices. Staging is $\repeatop$ on the storage, which exercises the
swizzled-repeat rule, together with a stage axis on the thread maps, and
slicing is $\restrict$ on the composites. Four properties are asserted address
by address. Each lane's global vector is contiguous and consecutive lanes in a
quad read consecutive 16-byte vectors. Every shared-memory vector is 16-byte
aligned and contiguous, so the swizzle does not break vectorization. The writer
and reader thread--value maps are each bijective onto all 512 staged elements.
And for every logical element, the byte the copy writes is the byte the mma
fetch reads, across two pipelines derived independently. No composite in the
suite is written by hand.

\section{Related work}\label{sec:related}

\textbf{CuTe/CUTLASS}~\cite{cutedocs,cecka} is the system under study; its
operations and their conditions are summarized in
Section~\ref{sec:intro}. \textbf{Shah}~\cite{shah} formalizes the algebra and gives the admissibility
conditions. They answer when two layouts fuse into one, which is a question
this algebra never asks. \textbf{Carlisle et al.}~\cite{categorical} give CuTe's algebra
categorical foundations. \textbf{Triton's linear layouts}~\cite{triton} model
layouts as $\mathbb{F}_2$-linear maps on address bits, gaining generic layout
conversion; that work is about what a
representation can express, where ours is about which maps may be built, and
the two combine: an $\mathbb{F}_2$ matrix is one way to implement a swizzle
here. \textbf{Bhaskaracharya et al.}~\cite{isl} model both CuTe and Triton layouts
as integer set relations and reason about composition, inversion, and
complement with a general Presburger solver; the question of
Section~\ref{sec:decide} is expressible there, and the point of our procedure
is that on a finite box it needs no solver---one pass of first differences per
axis. \textbf{Axe}~\cite{axe} unifies tiling, sharding, replication, and offsets
over named axes, from device meshes down to the lane, warp, and register
mappings of tensor-core tiles---the same grain as ours, over a wider span---and
makes replication explicit as we do. Its Appendix~A.2.1 contains the
characterization Section~\ref{sec:decide} builds on, with the attribution and
the difference stated there.
\textbf{Halide}~\cite{halide}, \textbf{TVM}~\cite{tvm}, and
\textbf{MLIR}~\cite{mlir} carry bounds-aware integer expression simplifiers,
which is the tool one reaches for first here, and
Section~\ref{sec:emit} reports why a decision procedure replaced ours.

Taken one at a time the distinctions this design rests on are familiar.
Replication is explicit in Triton and in Axe, the read and write asymmetry is
known to every kernel author, and slicing a view without rewriting it is
ordinary practice. What is new is that all three are enforced together:
replication is a constructor that may be read and not stored through, nothing
maps out of an address, and a restriction sits beside an unmodified layout and
so applies to composites. Those requirements, rather than any condition about
fusing, are what make a composition legal.

\section{Scope and limitations}\label{sec:scope}

The domain is dense tiled tensor-core kernels, the workload CuTe exists for,
and every claim here is scoped to it. Ragged tiles are handled by masking,
which is outside the algebra. Negative strides are unsupported. Tilers for
$\divideop$ that are layouts rather than shapes are deferred until a kernel
needs them. Density is decided by Lemma~\ref{lem:dense} without enumeration,
and the test suite checks that criterion against an enumeration oracle on 400
random cases; alias-freedom, injectivity, and anything about a composite cost
$O(\size)$ enumeration when the kernel is compiled, and the decision procedure
of Section~\ref{sec:decide} costs a further $O(\size)$ evaluations. Deciding
that question with extents known only at run time is outside this scope, where
tile shapes are compile-time constants.

\section{The language this serves}\label{sec:coda}

This algebra describes data movement for a warp-level GPU language we are
building, in which the unit of execution in the source is the warp, lanes are
SIMD data and divergence cannot be written, warps communicate through pipes
that compile to barrier machinery, and every buffer through a pipe carries a
layout from this paper. The language is future work and the algebra stands
without it.

\begin{thebibliography}{9}

\bibitem{cutedocs}
NVIDIA. \emph{CuTe: Layouts and Layout Algebra}. CUTLASS documentation,
\texttt{docs.nvidia.com/cutlass}, accessed 2026.

\bibitem{cecka}
C.~Cecka. \emph{CuTe Layout Representation and Algebra}. NVIDIA Research;
arXiv:2603.02298, 2026.

\bibitem{shah}
J.~Shah. \emph{A Note on the Algebra of CuTe Layouts}. Colfax Research
technical note, 2024.

\bibitem{categorical}
J.~Carlisle, J.~Shah, R.~Stern, P.~VanKoughnett.
\emph{Categorical Foundations for CuTe Layouts}. arXiv:2601.05972, 2026.

\bibitem{triton}
K.~Zhou, M.~Lezcano, A.~Goucher, A.~Rakhmati, J.~Niu, J.~Lebar, P.~Szczerbuk,
P.~Bell, P.~Tillet, T.~Raoux, Z.~Moudallal. \emph{Linear Layouts: Robust Code
Generation of Efficient Tensor Computation Using $\mathbb{F}_2$}. ASPLOS 2026;
arXiv:2505.23819.

\bibitem{isl}
S.~G.~Bhaskaracharya, A.~Acharya, B.~Hagedorn, V.~Grover. \emph{Modeling
Layout Abstractions Using Integer Set Relations}. arXiv:2511.10374, 2025.

\bibitem{axe}
B.~Hou, H.~Jin, G.~Wang, J.~Chen, Y.~Cai, L.~Yang, Z.~Ye, Y.~Ding, R.~Lai,
T.~Chen. \emph{Axe: A Simple Unified Layout Abstraction for Machine Learning
Compilers}. arXiv:2601.19092, 2026.

\bibitem{ptx}
NVIDIA. \emph{Parallel Thread Execution ISA}. Matrix fragment layout tables,
accessed 2026.

\bibitem{halide}
J.~Ragan-Kelley, C.~Barnes, A.~Adams, S.~Paris, F.~Durand, S.~Amarasinghe.
\emph{Halide: A Language and Compiler for Optimizing Parallelism, Locality,
and Recomputation in Image Processing Pipelines}. PLDI 2013.

\bibitem{tvm}
T.~Chen et al. \emph{TVM: An Automated End-to-End Optimizing Compiler for
Deep Learning}. OSDI 2018.

\bibitem{mlir}
C.~Lattner et al. \emph{MLIR: Scaling Compiler Infrastructure for Domain
Specific Computation}. CGO 2021.

\end{thebibliography}

\end{document}
